\section{Finite symmetry and the periodic reconstruction branches}
\label{sec:v49-finite-symmetry}

The local signed inverse and analytic image recovery do not require
asymmetry.  We separate the finite alternatives caused by rotational
symmetry from the continuous ambiguity present for a circle.  The same
distinction will sharpen the differential theorem in the next section.
Throughout this section the selected design is fixed and realizable.
All competing realizations use the same signed contact-law model on
an admissible collar at the prescribed offsets.

\subsection{Congruence choices and a local angular lift}

For a convex body $K$ let
\[
 s(K)=\frac1\pi\int_0^{2\pi}h_K(\theta)n(\theta)\,d\theta,
 \qquad
 z_k(K)=\frac1{2\pi}\int_0^{2\pi}h_K(\theta)e^{-ik\theta}\,d\theta.
\]
The center $s$ is proper-Euclidean equivariant, and $z_k$, $k\ge2$,
is invariant under translations.  Write $\operatorname{Sym}^+(K)$
for the group of proper Euclidean isometries preserving $K$.

\begin{lemma}[Finite congruence alternatives]
\label{lem:v49-finite-congruences}
If $K$ is noncircular, $\operatorname{Sym}^+(K)$ is a finite cyclic
group of order
\[
                m(K)=\gcd\{k\ge2:z_k(K)\ne0\}.
\]
For two properly congruent copies $C_s,C_t$ of $K$, the set of proper
congruences between them has exactly $m(K)$ members.  Their translations
are determined by their rotations and the two centers.
\end{lemma}
\begin{proof}
Every symmetry fixes $s(K)$.  After centering it is a rotation.
A centered body is a disk precisely when all harmonics of orders
$k\ge2$ vanish: Fourier uniqueness then makes its support function
constant.  Otherwise choose one nonzero $z_k$.  An allowed rotation
$e^{i\alpha}$ must satisfy $e^{-ik\alpha}=1$, so there are finitely
many such rotations.  A rotation preserves the support function
precisely when this equality holds for every nonzero harmonic.
The common solutions are the $m(K)$th roots of unity.  The gcd of the
infinite nonempty set is already the gcd of a finite subset, since
successive positive integer gcds eventually stabilize.  Given one
congruence, all others differ by a symmetry of $K$.  For each selected
rotation $U$, its translation is $s(C_t)-Us(C_s)$.
\end{proof}

\begin{lemma}[Local alignment without a global phase choice]
\label{lem:v49-local-alignment}
Let $C_s(t),C_t(t)$ be $C^1$ families of properly congruent support
images, with noncircular base images.  Suppose their actual proper congruences
form a continuous family through the base alignment $U_0=e^{i\alpha_0}$.
Then one nonzero base harmonic $z_k$,
$k\ge2$, determines a $C^1$ local branch of their rotation and
translation through that alignment.  If both image derivatives vanish,
the derivatives of that rotation and translation vanish.  The assertion
does not require the symmetry group to be constant along the family.
\end{lemma}
\begin{proof}
Choose $k$ with $z_k(C_s(0))\ne0$; the corresponding coefficients in
both copies remain nonzero.  Put $r(t)=z_k(C_t(t))/z_k(C_s(t))$ and
$r_0=e^{-ik\alpha_0}$.  In a small disk about $1$ use the logarithm
with $\operatorname{Log}1=0$.  Along the actual congruent family,
\begin{equation}\label{eq:v49-local-phase}
 U(t)=U_0\exp\{-k^{-1}\operatorname{Log}(r(t)/r_0)\},\qquad
 a(t)=s(C_t(t))-U(t)s(C_s(t)).
\end{equation}
Indeed, choose the continuous actual angular lift through $\alpha_0$
and restrict it by $|\alpha(t)-\alpha_0|<\pi/k$.  Then
$\operatorname{Log}(r(t)/r_0)=-ik(\alpha(t)-\alpha_0)$, proving
\eqref{eq:v49-local-phase}.  The displayed formula is $C^1$, and
\begin{equation}\label{eq:v49-angular-velocity}
             \frac{\dot r}{r}=-ik\dot\alpha,
                 \qquad \dot\alpha=\frac{i}{k}\frac{\dot r}{r}.
\end{equation}
The last expression is real along the congruent family.  Stationary
support images have zero center and harmonic derivatives, so
$\dot\alpha=0$ and $\dot a=0$.  Other congruence branches can exist,
but are separated from the selected angular lift.  The branch followed by the actual family is the one required here;
a base symmetry which is broken need not extend to nearby pairs.
Symmetry breaking does not affect the nonzero harmonic or this argument.  No phase is selected globally from unlabelled
noisy copies by \eqref{eq:v49-local-phase}.
\end{proof}

\subsection{A finite reconstruction of the exact fiber}

Let $E=\mathcal T\cup\{f_1,f_2\}$ be a clear skeleton as in
Theorem~\ref{thm:v45-clear-skeleton}.  Its graph on obstacle labels is
connected, and its two independent cycle gains form a fixed invertible
integer matrix $M$.  In addition, form a graph whose vertices are the
channels in $E$, connecting two channels whenever they contain an
occurrence of the same obstacle label.  This graph is connected: a
path in the obstacle graph gives a chain of incident channels.
Choose a rooted spanning tree $\mathcal U$ in this channel graph.
For each link $\nu$ of $\mathcal U$ fix the particular two occurrences
of its common obstacle label, denoted $a(\nu)$.  These are combinatorial
choices, including for loops and multiple edges, not supplied poses.

\begin{theorem}[Finite fibers for noncircular obstacles]
\label{thm:v49-finite-fiber}
Suppose every obstacle of a realizable analytic table $T_0$ is
noncircular.  Its gaps and same-type signed laws at one known sufficiently
small positive offset per channel determine all realizations of the
same marked datum, modulo a common proper Euclidean motion, as a finite
set.  With the tree $\mathcal U$ just chosen its cardinality is at most
\begin{equation}\label{eq:v49-fiber-bound}
                         \prod_{\nu\in\mathcal U}m(K_{a(\nu)}).
\end{equation}
The bound counts congruence assignments and need not be attained.
When every $m(K_a)=1$, the fiber consists of a single table.
\end{theorem}
\begin{proof}
Theorem~\ref{thm:v26-density-inverse}, the smooth signed jet inverse,
and analytic continuation recover the entire pair of obstacle images
$C_{e,-},C_{e,+}$ in every channel frame.  This step uses no proper
asymmetry.  It also determines their centers and finite symmetry orders.
Every realization of the datum must have these same images.

Fix the rotation of the root channel as the identity.  Propagate
rotations $R_e$ along $\mathcal U$.  At a link matching occurrences
$C_s$ and $C_t$, Lemma~\ref{lem:v49-finite-congruences} supplies exactly
$m(K_{a(\nu)})$ proper congruences $C_t=UC_s+a$.  For each choice set
$R_t=R_sU^{-1}$.  Thus $R_t(C_t-s(C_t))=R_s(C_s-s(C_s))$.
All assignments are finite and their number is bounded by
\eqref{eq:v49-fiber-bound}.  Their translations are not propagated as
unknown lattice displacements; only the rotations have been fixed so far.

For a complete assignment, require all centered, rotated occurrences
of each obstacle to agree, not merely those on $\mathcal U$.
Assignments failing this test are discarded.  Denote the resulting
centered representative images by $\widehat K_a$.  Orient the obstacle
tree $\mathcal T$ from a root obstacle and define
\begin{equation}\label{eq:v49-candidate-cochain}
 \begin{split}
 D_e&=R_e\{s(C_{e,+})-s(C_{e,-})\},\qquad p_1=0,\\
 p_a&=\sum_{e\in[1,a]_{\mathcal T}}\epsilon_e D_e,\qquad
 m_a=\sum_{e\in[1,a]_{\mathcal T}}\epsilon_e\ell_e,\\
 v_i&=p_{a_i}+D_{f_i}-p_{b_i},\qquad
 \eta_i=m_{a_i}+\ell_{f_i}-m_{b_i}\quad(i=1,2),\\
 L&=(v_1\ v_2)M^{-1},\qquad M=(\eta_1\ \eta_2),\\
 c_a&=p_a-Lm_a,\qquad K_a=\widehat K_a+c_a.
 \end{split}
\end{equation}
This uses the actual $M^{-1}$, without assuming $\det M=1$.
The cochain identities give, for every tree edge and both remaining
edges,
\begin{equation}\label{eq:v49-edge-consistency}
                 c_b+L\ell_e-c_a=D_e.
\end{equation}
For tree edges this follows directly by telescoping; for $f_i$ it is
$L\eta_i=v_i$.  Each assignment therefore produces at most one marked
candidate in the root-center and root-direction gauge.

Retain a candidate only when $\det L>0$, the lifted obstacles are
pairwise disjoint, and the selected pairs have the stipulated clear
contacts and collars.  These conditions are part of realization, not
consequences of the integer gain matrix.  They require only finitely
many lattice checks once a nonsingular candidate has been produced:
if all representative bodies and selected segments lie in a ball of
radius $R$, a translated object can meet one of them only for marks
bounded by a constant times $R/s_{\min}(L)$.  Exact shape equality is
an equality of the already known analytic images.  This description is
a finite-branch inverse of exact function-valued data, not a numerical
decision algorithm for arbitrary finite precision inputs.

Every actual realization gives one of the congruence assignments and
satisfies \eqref{eq:v49-candidate-cochain}, hence occurs among the
retained candidates.  Conversely, for a retained candidate, centered
shape agreement and \eqref{eq:v49-edge-consistency} place each entire
ordered pair exactly as its recovered channel pair, by one proper
motion.  Its closest segment, oriented transverse convention and local
actions and amplitudes therefore agree with those recovered from the
datum.  The relative law theorem yields exactly the same normalized
conditional laws.  Global phase-volume normalization cancels under
conditioning; it is not extra input here.  All selected design
conditions have been tested.  Thus every retained candidate is a
realization.  Duplicate candidates can be identified in the common
gauge.  Finally, realizability makes the set nonempty, and when all
symmetry orders are one the upper bound is one.
\end{proof}

The theorem removes no discrete ambiguity by an arbitrary phase rule.
Nor does it assert global uniqueness for all noncircular tables.
Its differential counterpart uses the local branch through the actual
base alignment and therefore needs only the absence of a continuous
rotational stabilizer.

\subsection{A circular comparison for the same selected design}

\begin{proposition}[A circle and a variable marked Gram form]
\label{prop:v49-circular-comparison}
For one disk orbit of radius $1/10$, let
\[
 L_t=\begin{pmatrix}1&-\sin t\\0&\cos t\end{pmatrix},
                         \qquad |t|<1/10.
\]
The two selected channels with gains $e_1,e_2$ are clear on a common
neighborhood and have constant gaps and same-type signed limiting
conditional laws at a common sufficiently small positive offset.
Nevertheless the marked lattice Gram form has a nonzero derivative
at zero, and the variation is not a common Euclidean motion.
\end{proposition}
\begin{proof}
Both columns of $L_t$ have length one.  Thus the two disk pairs have
gap $4/5$ and each ordered pair is constant up to its own proper
channel motion.  The same-type local actions and amplitudes, and hence
their conditional laws, are constant.  Their common offset and collars
can be chosen after clearance is verified.

Put $\epsilon=\|L_t-I\|=2|\sin(t/2)|<1/10$.
For every nonzero integer mark $m$, $|L_tm|\ge(1-\epsilon)|m|$,
so distinct disk centers are separated by at least $9/10$.
Consider either unit center segment from $0$ to $L_te_i$.
If a third center $L_tm$ had distance at most $1/2$ from that segment,
then $|L_tm|\le3/2$, whence $|m|\le5/3$.  Apart from the two incident
centers, $m$ is therefore one of the eight nearest nonzero integer
neighbors.  At $t=0$ its distance from $[0,e_i]$ is at least one.
Perturbing that center and segment decreases the distance by at most
$(|m|+1)\epsilon\le(\sqrt2+1)/10<1/4$.
This contradicts the assumed distance at most $1/2$.
Consequently every third disk has distance greater than $2/5$ from
the center segment, and also from its shorter boundary-to-boundary
closest segment.  Uniformly small gates are admissible.

Finally,
\[
 L_t^TL_t=\begin{pmatrix}1&-\sin t\\-\sin t&1\end{pmatrix},
 \qquad
 \dot L_0=\begin{pmatrix}0&-1\\0&0\end{pmatrix},\qquad
 \left.\frac{d}{dt}\right|_0 L_t^TL_t=
                \begin{pmatrix}0&-1\\-1&0\end{pmatrix}.
\]
A common rotation would preserve the Gram form.  Thus this is not
such a motion.  The example concerns the two specified channels;
no finite-horizon condition is asserted.
\end{proof}

This comparison is outside the noncircular class and does not contradict
the exact or differential theorem.  It exhibits the geometric role of
the hypothesis for this design, not a universal failure of reconstruction
from other designs or richer observations.
